August 3, 2026

Editor-in-Chief
Computers & Security
Elsevier

SUBJECT: Submission of Manuscript P7 — Bounded Approximate Nearest Neighbor Identity Retrieval in Privacy-Constrained Environments

Dear Editor-in-Chief and Associate Editors,

We are pleased to submit our original research manuscript titled "Bounded Approximate Nearest Neighbor Identity Retrieval in Privacy-Constrained Environments" for publication consideration in Computers & Security (Elsevier).

1. SUBSYSTEM BACKGROUND AND MOTIVATION
High-scale biometric face recognition in open-set campus settings must handle thousands of enrolled students while rejecting unenrolled visitors. Conventional closed-set facial matching fails under open-set conditions, and linear vector searches scale poorly. This paper presents an open-set identity retrieval architecture that guarantees sub-millisecond search performance without persistent biometric image storage.

2. CORE TECHNICAL CONTRIBUTIONS AND NOVELTY
  - Open-Set ArcFace & FAISS Indexing: Combines an ArcFace additive angular margin loss (512-dimensional vector space) with an Inverted File (IVF-PQ) FAISS index, achieving 99.2% Open-Set Identification Rate (OSIR) and 99.5% Unknown Identity Rejection Rate (UIRR).
  - Scalability to 100,000 Identities: Benchmarks vector retrieval across galleries scaling up to 100,000 student profiles, sustaining sub-millisecond query search latencies (0.8ms) within a 200MB memory footprint.
  - Volatile Vector Boundary: Operates strictly on normalized vector embeddings, enforcing zero-disk retention of biometric face crops.

3. SUITABILITY FOR COMPUTERS & SECURITY
This paper directly aligns with Computers & Security's focus on biometric security systems, high-performance vector retrieval, privacy protection mechanisms, and scalable authentication infrastructure.

4. ETHICS AND ORIGINALITY DECLARATIONS
This manuscript is an original work, has not been published previously, and is not currently submitted to any other publication venue.

Sincerely,

Polisetti Narendra (Corresponding Author)
Department of Computer Science & Engineering
Swarnandhra College of Engineering and Technology (Autonomous)
Seetharampuram, Narasapur-534280, Andhra Pradesh, India
Email: narendraa1996@gmail.com
